The USB-C connector acted as a sink of of power flow \cite{ti-usb-typec-guide}. The schematic was mostly copied over from v0. I selected the \href{https://au.mouser.com/ProductDetail/Same-Sky/UJ20-C-H-G-MSMT-2-P16-TR?qs=IKkN%2F947nfBldcIwp9d6BA%3D%3D}{UJ20-C-H-G-MSMT-2-P16-TR} component over top-mounts because FISH and the ADCS Integration Board had their connectors being torn off, the mid-mount provided more security to the board and is used by the EPIC board. I was able to remove the power sequencer (since the U5 does not require a specific startup sequence) and replace the buck converter with an LDO (because of the lower maximum power consumption). In addition, I removed the differential data pair (D+, D-) because Robert Howie suggested the COM port functionality was not very useful compared to using the ST-Link.
\nl
As outlined in pg. 54 of \cite{ti-usb-typec-guide}, the sink required 5.1k$\Omega$ Rd resistors on the two CC lines. Since the total capacitance of VBUS did not exceed 10$\mu$F, I did not need a power path or load switch. Moreover, I was not going to exceed 15W of power consumption, so a USB PD controller was not required.
\nl
Using a Current Sense Amplifier (CSA), I estimated the power consumed by the circuit by placing a low-value shunt resistor before the LDO and measuring the voltage drop across them from V1 to V2 wth a CSA. Any voltage losses created by the resistor would be resolved by the LDO regulator to +3V3. The power consumption data would infer what the maximum current carrying capacitor of traces are for the prototype; and it will provide the Systems Engineering Lead with vital power requirements for the CubeSat. With some miscalculations, I found that the resistor needed to be rated for 3W with an impedance of 2m$\Omega$.
\nl
The bi-directional TVS diode acted as surge protection by clamping the voltage in both directions (copied over from ADCS Integration Board). For more information about the pi-filter network, please see \autoref{sec:appendices}.

\paragraph{Remark}
The maximum ADC is when the CSA outputs 3.3 V \cite{TI-Manchukonda2019ShuntResistorQuiz}. The INA181A2 has a gain of 50 according to the \href{https://www.ti.com/lit/ds/symlink/ina181.pdf?ts=1775027611016&ref_url=https%253A%252F%252Fau.mouser.com%252F}{datasheet}, which means the maximum voltage difference across the shunt must be $3.3/50=66$mV. For a maximum current of 1 A, the maximum impedance must be $66\mu\Omega$. The \href{https://au.mouser.com/ProductDetail/TE-Connectivity-Holsworthy/TLRP3A30CR068FTE?qs=QNEnbhJQKvaIeoPRm6mvGw%3D%3D}{selected replacement} had a resistance of $68\mu\Omega$ as $66\mu\Omega$ is not a common resistor value.
\nl
The transplant process was not easy because of the component's large size. The large cross sectional area meant it was not possible to cut the component in half and remove each pad independently from each other and the large pad sizes meant it was difficult to heat the solder since the heat quickly dissipated into the PCB.
With help from Aman Tanday, he was able carefully heat one pad and lift the same side to remove the resistor.
\nl
The installation process was straight-forward because the component was so large. Some flux and solder did the job. After the resistor was placed, we had no issues powering the board.